\documentclass[11pt]{article}
\usepackage{fontspec,amsmath,amssymb,amsthm,hyperref}
\usepackage[margin=1in]{geometry}
\setlength{\emergencystretch}{3em}
\title{The first-order scalar accumulation and its exact scope}
\author{LangevinLean mathematical exposition}
\date{}
\begin{document}\maketitle
\noindent Source module: \texttt{KineticLangevin/Recurrence.lean}. Compile this proof document with XeLaTeX or LuaLaTeX. The accompanying HTML pairs the mathematical argument with the exact Lean statements and records the scope of the formalization.


The numerical accumulation part of the KLMC proof begins with an explicit one-step coupling inequality. This module derives the stable contraction coefficient and sums its affine recurrence, retaining the coupling inequality and reference velocity norm as visible hypotheses. It therefore checks the numerical implication without claiming the still-open Gibbs stationarity and actual Gaussian-chain Wasserstein endpoint.


\section{Prove the one-sided exponential quadratic estimate}


For \(x\ge0\), let \(F(x)=1-x+x^2/2-e^{-x}\). Its derivative is \(-1+x+e^{-x}\), nonnegative because \(e^{-x}\ge1-x\). Therefore \(F\) is nondecreasing and \(F(x)\ge F(0)=0\), proving \(e^{-x}\le1-x+x^2/2\). Lean constructs the derivative through the polynomial and exponential derivative rules, proves its sign with \texttt{Real.add\_one\_le\_exp}, applies the derivative monotonicity criterion, and evaluates the zero endpoint. The estimate is proved analytically rather than introduced as an unchecked Taylor remainder.



\noindent Lean declarations: \texttt{exp\_neg\_le\_quadratic}.

\section{Keep most of the exact-flow contraction after the local correction}


Set \(r=m/\gamma\). Suppose \(r,M,h\ge0\), \(2r^2\le M\), and \(4Mh\le r\). The quadratic exponential estimate gives \(e^{-rh}+3Mh^2/4\le1-rh+r^2h^2/2+3Mh^2/4\). Since \(r^2/2\le M/4\), the positive correction is at most \(Mh^2\le rh/4\). Thus the total is at most \(1-3rh/4\). Lean multiplies the curvature and step inequalities by the nonnegative time powers and combines them with \texttt{nlinarith}. The next theorem derives both scalar conditions from \(m>0\), \(m\le M\), \(\gamma^2\ge m+M\) and \(4\gamma Mh\le m\), dividing only by positive \(\gamma\) or \(\gamma^2\).



\noindent Lean declarations: \texttt{firstOrder\_coefficient\_bound}, \texttt{kinetic\_scalar\_conditions}.

\section{Sum the affine recurrence with the visible local-coupling premise}


Assume nonnegative errors satisfy \(e_{n+1}\le e^{-rh}e_n+3Mh^2(C+e_n)/4\), with \(r,h>0\), \(C\ge0\), and the scalar restrictions above. Put \(q=1-3rh/4\). The coefficient bound and positivity of the exponential prove \(0\le q<1\), and multiplying that bound by \(e_n\ge0\) yields \(e_{n+1}\le qe_n+3Mh^2C/4\). The shared affine-recurrence theorem gives \(e_n\le q^ne_0+(3Mh^2C/4)/(1-q)\). Positive \(r,h\) permit denominator cancellation, leaving \(e_n\le q^ne_0+MhC/r\). Lean uses \texttt{field\_simp} and \texttt{ring} only for that final quotient. The one-step coupling inequality is an explicit assumption, so this result must not be presented as the full actual-law KLMC theorem.



\noindent Lean declarations: \texttt{firstOrder\_coupling\_recurrence\_bound}.

\section{Restore the paper’s physical parameters and position coefficient}


Substitute \(r=m/\gamma\) into the scalar recurrence theorem. The derived scalar conditions justify the call, and rational simplification gives \(e_n\le(1-3mh/(4\gamma))^ne_0+\gamma MhC/m\). Suppose in addition that \(d_n\le\sqrt2e_n/\gamma\) and \(e_0\le\gamma d_0\). The nonnegative projection factor and nonnegative power of the proved stable contraction coefficient allow multiplication of these inequalities. Canceling \(\gamma\) yields

\[d_n\le\sqrt2(1-3mh/(4\gamma))^nd_0+Mh\sqrt2C/m.\]

Lean explicitly proves the projection factor and contraction coefficient nonnegative before transporting inequalities, then uses \texttt{field\_simp} for the scalar cancellations. Setting \(C=\sqrt p\) matches the paper’s numerical expression. The local coupling, projection and initialization bounds remain explicit hypotheses on sequences; neither declaration identifies \(d_n\) with actual Wasserstein error or constructs a stationary reference.



\noindent Lean declarations: \texttt{kinetic\_firstOrder\_coupling\_bound}, \texttt{kinetic\_firstOrder\_projected\_bound}.

\end{document}
